\scanpage{56}
\mspage{55}

Je rappelle pour mémoire les conditions \inserted{Sous at 1,2} \inserted{de page 42} pour qu'une partie ${\unsure{\mathcal F|\mathscr F}}^{\unsure{f|*|\prime}}$ de $\operatorname{Fig}(L)$ (en l'occurrence) :

\uline{At ens$^{*}$} 2 \quad $\mathcal F$ est fermé dans $\operatorname{Fig}(L)$ pour $\leq$, i.e. si $F\in\mathcal F$, alors \unsure{toute|tt} sous figure de $F$ est dans $F$.

\uline{At ens$^{*}$} 3 \quad Si $F,G\in\mathcal F$, et si pour \deleted{pour \unsure{tte} $F$ et $G$ sont compat} toute multistrate $F_A$ de $F$, $G_B$ de $G$, $\{F_A\unsure{,|\vee}G_B\}$ est majoré dans $\mathcal F$, \inserted{i.e. $F_A\mathrel{\vcenter{\hbox{\begin{tikzpicture}[x=1ex,y=1ex,baseline]\draw (0,0)--(1.05,.50)--(0,1);\draw (1.05,0)--(0,.50)--(1.05,1);\end{tikzpicture}}}}G_A$ et $F_A\unsure{\dot\vee|\vee}G_A\in\mathcal F$} \deleted{alors} \inserted{(a fortiori les figures \deleted{\unsure{ou}} $F,G$ sont \inserted{alors} \unsure{ensemblistmt} compatibles)} \deleted{$F,G$} $F\vee G\in\mathcal F$.

Ceci peut s'expliciter \unsure{aussi|ainsi} \inserted{(partiellement)} en termes d'un ensemble $M$ de figures \inserted{(ensemblistes)} élémentaires dans $L$, \deleted{\unsure{en convenant pour} $\mathcal F$} \inserted{avec une} relation $\underset{M}{\mathrel{\vcenter{\hbox{\begin{tikzpicture}[x=1ex,y=1ex,baseline]\draw (0,0)--(1.05,.50)--(0,1);\draw (1.05,0)--(0,.50)--(1.05,1);\end{tikzpicture}}}}}$ (``compatibilité géométrique'') \inserted{réfl. et sym} dans $M$, satisfaisant les conditions

\[
\big|\,\text{\uline{At ens}}\ \unsure{\sharp|\mathrm{gf}|{}^{**}}\ 1\qquad\forall\,x\in L,\quad\{\{x\}\}\in M
\]

\noindent\raisebox{0pt}[0pt][0pt]{\rule[-1.5\baselineskip]{.35pt}{2.4\baselineskip}}\quad
\uline{At ens} $\sharp$ 2 \quad \deleted{\unsure{$A$}} $\forall\,X\in M$, toute sous-multistrate est dans $M$

\noindent\raisebox{0pt}[0pt][0pt]{\rule[-2.1\baselineskip]{.35pt}{3\baselineskip}}\quad
\uline{At ens} $\deleted{{}^{**}}\sharp$ 3 \quad Si $X,Y\in M$ et $X\underset{M}{\mathrel{\vcenter{\hbox{\begin{tikzpicture}[x=1ex,y=1ex,baseline]\draw (0,0)--(1.05,.50)--(0,1);\draw (1.05,0)--(0,.50)--(1.05,1);\end{tikzpicture}}}}}Y$, et $X',\deleted{\unsure{Y'}}\leq X$, $Y'\leq Y$, alors $X'\underset{M}{\mathrel{\vcenter{\hbox{\begin{tikzpicture}[x=1ex,y=1ex,baseline]\draw (0,0)--(1.05,.50)--(0,1);\draw (1.05,0)--(0,.50)--(1.05,1);\end{tikzpicture}}}}}Y'$ ; de plus $X\underset{M}{\mathrel{\vcenter{\hbox{\begin{tikzpicture}[x=1ex,y=1ex,baseline]\draw (0,0)--(1.05,.50)--(0,1);\draw (1.05,0)--(0,.50)--(1.05,1);\end{tikzpicture}}}}}Y\Longrightarrow X\underset{\mathrm{ens}}{\mathrel{\vcenter{\hbox{\begin{tikzpicture}[x=1ex,y=1ex,baseline]\draw (0,0)--(1.05,.50)--(0,1);\draw (1.05,0)--(0,.50)--(1.05,1);\end{tikzpicture}}}}}Y$
